\chapter{样式使用说明}
\label{cha:command}


\section{参考文献}
\textcolor{red}{博士学位申请人的文献阅读量一般不少于100篇，其中外文文献一般不少于1/3；近五年的论文一般不少于1/3；绪论部分应对所读文献加以分析和综合。}

参考文献格式

中文书刊：作者按中文写法，姓在前、名在后；英文书刊：作者按英文习惯写法，如名在前、姓在后，名用首字母缩写、姓用全称。一般6人以内须列出全部作者，6人以上写6人再加“等”（英文加“et al”））。每个参考文献的最后不加标点符号。

（1）图书：最多列出6个作者，作者与作者之间用逗号分隔. 书名. 版本（第×版）. 译者. 出版地: 出版者, 出版年. 起页-止页（可选）

（2）期刊：最多列出6个作者，作者与作者之间用逗号分隔. 文章名. 期刊名（全称）. 年号, 卷号（期号）: 起页-止页或论文编号

（3）会议论文集：最多列出6个作者，作者之间用逗号分隔. 文章名. 见（英文用“in”）：会议名称（或论文集）. 会议城市, 国家, 会议时间, 出版者, 出版年: 起页-止页

（4）专利：专利申请者. 专利题名. 专利国别, 专利文献种类, 专利号, 出版年

（5）学位论文：作者. 题名：[博士（或硕士）学位论文]. 保存地点: 保存单位（如华中科技大学, 年份）

参考文献（举例）
\begin{enumerate}
\renewcommand{\labelenumi}{[\theenumi]}
\item	闫明礼, 张东刚. CFG桩复合地基技术及工程实践（第二版）. 北京: 中国水利水电出版社, 2006

\item	M. Chalfie, S. R. Kain. Green fluorescent protein: properties, applications, and protocols. Hoboken, New Jersey: Wiley-interscience, 1998

\item	詹向红, 李德新. 中医药防治阿尔茨海默病实验研究述要. 中华中医药学刊, 2004, 22(11): 2094-2096

\item	E. S. Lein, M. J. Hawrylycz, N. Ao, M. Ayres, A. Bensinger, A. Bernard, et al. Genome-wide atlas of gene expression in the adult mouse brain. Nature, 2007, 445(7124): 168-176

\item	M. L. Bouxsein, S. K. Boyd, B. A. Christiansen, R. E. Guldberg, K. J. Jepsen, R. Müller. Guidelines for assessment of bone microstructure in rodents using micro–computed tomography. Journal of Bone and Mineral Research, 2010, 25(7): 1468-1486

\item	Y. Shunsuke, A. Masahide, K. Masayuki, M. Yoshizawa. Performance evaluation of phase-only correlation functions from the viewpoint of correlation Filters, in: 2018 Asia-Pacific Signal and Information Processing Association Annual Summit and Conference (APSIPA ASC), Honolulu, HI, USA, 12-15 Nov. 2018, Proceedings of the IEEE, 2019: 1361-1364

\item	T. Yao, J. Wan, P. Huang, X. He, F. Wu, C. Xie. Building efficient key-value stores via a lightweight compaction tree. ACM Transactions on Storage, 2017, 13(4): 1-28

\item	刘加林. 多功能一次性压舌板: 中国, ZL92214985. 2 [P]. 1993

\item	李清泉. 基于混合数据结构的三维GIS数据模型与空间分析研究[博士学位论文]. 武汉: 武汉测绘科技大学, 1998

\end{enumerate}

\textcolor{red}{以下是几个例子, 注意, 中文参考文献需要额外增加一个 lang="chinese"}
~\\
期刊论文举例：

Pan等人提出了一种全新的算法\cite{pan2018classification, pan2014cell, song2013asynchronous, lin2022adaptive, liu2023low}。

~\\
会议论文举例：

XXXXxx\cite{li2021large}

~\\
图书举例：

xxxxxx\cite{shavlik1990readings, yan2006CFG}





